\documentclass[11pt,a4paper]{article}

\usepackage[utf8]{inputenc}
\usepackage[T1]{fontenc}
\usepackage{lmodern}
\usepackage[margin=0.78in]{geometry}
\usepackage[table]{xcolor}
\usepackage{array}
\usepackage{tabularx}
\usepackage{booktabs}
\usepackage{enumitem}
\usepackage{titlesec}
\usepackage{fancyhdr}
\usepackage{lastpage}
\usepackage{hyperref}

\definecolor{TextInk}{HTML}{233142}
\definecolor{HeadingBlue}{HTML}{1D3557}
\definecolor{AccentTeal}{HTML}{2A7F92}
\definecolor{AccentGold}{HTML}{C58B4E}
\definecolor{TableStripe}{HTML}{F3F8FA}
\definecolor{SoftBlue}{HTML}{EEF6F8}
\definecolor{SoftGreen}{HTML}{EEF7F0}
\definecolor{SoftGold}{HTML}{FCF6E8}

\hypersetup{
  colorlinks=true,
  linkcolor=HeadingBlue,
  urlcolor=AccentTeal,
  pdftitle={Geography Paper 1 SL 2023 Student Work Marking Report},
  pdfauthor={IB Geography SL Preparation}
}

\color{TextInk}
\setlength{\parskip}{0.45em}
\setlength{\parindent}{0pt}
\setlength{\tabcolsep}{5pt}
\renewcommand{\arraystretch}{1.18}
\setlist[itemize]{leftmargin=1.35em,itemsep=3pt,topsep=4pt}

\newcolumntype{Y}{>{\raggedright\arraybackslash}X}
\newcolumntype{L}[1]{>{\raggedright\arraybackslash}p{#1}}
\newcommand{\term}[1]{\textcolor{HeadingBlue}{\textbf{#1}}}
\newcommand{\score}[1]{\textcolor{AccentTeal}{\textbf{#1}}}
\newcommand{\boxnote}[3]{%
\noindent\colorbox{#1}{%
  \begin{minipage}{0.965\textwidth}
  \sffamily
  \textbf{\textcolor{HeadingBlue}{#2}} #3
  \end{minipage}
}}

\titleformat{\section}
  {\normalfont\Large\sffamily\bfseries\color{HeadingBlue}}
  {}{0pt}{}
  [\vspace{0.08em}{\color{AccentGold}\titlerule[0.8pt]}]

\titleformat{\subsection}
  {\normalfont\large\sffamily\bfseries\color{AccentTeal}}
  {}{0pt}{}

\pagestyle{fancy}
\setlength{\headheight}{14pt}
\fancyhf{}
\fancyhead[L]{\sffamily\color{AccentTeal}Teacher Marking Report}
\fancyhead[R]{\sffamily\color{HeadingBlue}May 2023 Paper 1}
\fancyfoot[C]{\sffamily\color{HeadingBlue}\thepage\ of \pageref{LastPage}}
\renewcommand{\headrulewidth}{0.4pt}

\begin{document}

\begin{center}
  {\sffamily\bfseries\color{HeadingBlue}\LARGE Geography Paper 1 SL 2023 Marking Report}\\[0.25em]
  {\sffamily\color{AccentTeal}\large Options F and G Student Work}
\end{center}

\boxnote{SoftBlue}{Basis for marking.}{This report marks
\path{plans/Geography_paper_1_SL_-2023-student-work.md} against the official
May 2023 HLSL Paper 1 markscheme in
\path{plans/Geography_paper_1__HLSL_markscheme-2023.pdf}. The student answered
Option F, Questions 11 and 12(b), and Option G, Questions 13 and 14(b).
Spelling and grammar were ignored unless they affected geographical meaning.}

\section{Overall Mark}

\rowcolors{2}{TableStripe}{white}
\begin{tabularx}{\textwidth}{L{0.16\textwidth}L{0.14\textwidth}Y}
\toprule
\textbf{Question} & \textbf{Mark} & \textbf{Teacher rationale} \\
\midrule
Q11 & \score{10/10} & Correct source answers for lettuce and cassava. The
mechanization answer earns credit when read as a reduction in human/manual
energy. Both food-insecurity explanations give clear resource or land-use
mechanisms and link these to increased food availability. \\
Q12(b) & \score{7/10} & Reaches the 7--8 band. The response is structured,
uses the Green Revolution/India as an applied example, and explains two
physical-geography influences: suitability of soil/climate and barriers such as
mountains, rivers and oceans. It remains at the lower end of the band because
the range of geographic factors is narrow and the case-study detail contains
some imprecision. \\
Q13 & \score{10/10} & All short-answer parts meet the markscheme: Kinshasa,
82, an environmental problem linked to cars/emissions, an economic
job-opportunity migration chain, and a demographic youthful-population/natural
increase chain. \\
Q14(b) & \score{8/10} & Reaches the 7--8 band securely. The Detroit example is
relevant and supported with population data, suburbanization/white flight,
industrial relocation, tax-base decline, infrastructure stress and social
tension. The answer has some balance and a judgment, but does not reach 9--10
because counter-urbanization, suburban residential impacts and stakeholder
perspectives are underdeveloped. \\
\midrule
\textbf{Total} & \score{35/40} & Strong two-option performance. The student is
secure on data-response questions and writes relevant 10-mark answers. The main
lost marks come from limited breadth and precision in the extended responses. \\
\bottomrule
\end{tabularx}
\rowcolors{2}{}{}

\section{Option F: Food and Health}

\subsection{Q11: Energy Inputs and Food Insecurity \score{10/10}}

\rowcolors{2}{TableStripe}{white}
\begin{tabularx}{\textwidth}{L{0.16\textwidth}L{0.12\textwidth}Y}
\toprule
\textbf{Part} & \textbf{Mark} & \textbf{Rationale} \\
\midrule
11(a)(i) & \score{1/1} & The student identifies \term{Lettuces}, matching the
lowest energy output on the graph. \\
11(a)(ii) & \score{1/1} & The student identifies \term{Cassava}, matching the
highest energy efficiency. \\
11(b) & \score{2/2} & The response states that machines require less energy
than manual methods for a task. This is creditworthy under the markscheme's
allowance for \term{less human energy}. For a safer answer, the student should
name \term{human/manual energy} explicitly, because total fossil-fuel energy
inputs may increase with mechanization. \\
11(c)(i) & \score{3/3} & Shows understanding of in vitro meat as meat not
requiring animal rearing/butchering, develops the reduced water, feed and land
requirements, and links this to more food being produced. This is a generous
full-credit reading because the definition is implicit; an explicit phrase such
as \term{lab-grown} or \term{synthetic meat} would be safer. \\
11(c)(ii) & \score{3/3} & Explains that vertical farming uses much less land,
can produce crops in urban areas with limited agricultural land, and can reduce
distribution needs while increasing food availability. This is a generous
full-credit reading because the definition is implicit; an explicit phrase such
as \term{crops grown in stacked layers} or \term{crops grown in buildings} would
be safer. \\
\bottomrule
\end{tabularx}
\rowcolors{2}{}{}

\subsection{Q12(b): Diffusion of Agricultural Innovation \score{7/10}}

The answer satisfies the 7--8 band because it is clearly focused on the
question and gives developed explanation rather than a list. It argues that
geographic setting affects whether an agricultural innovation is adopted, then
uses the Green Revolution/India example to show how soil and climate suitability
can encourage adoption. It also explains how mountains, rivers and oceans can
slow the human spread of ideas.

The mark is held at \score{7/10} because the range of factors is mostly
physical. The markscheme expects a broader understanding of geographic factors,
including economic factors such as capital and infrastructure, social factors
such as education and training, and political factors such as government action
or stability. The Green Revolution example is relevant, but the details are
imprecise: the answer attributes a rice seed in Mexico to Norman Borlaug, while
the stronger historical link would distinguish Borlaug's wheat work from later
HYV rice diffusion. The conclusion is relevant but not fully evaluative.

To move this to \score{9--10}, the student should compare several factors:
physical suitability, access to capital, infrastructure/communications,
government support, education/outreach and stakeholder power. The final
judgment should weigh which factors most affect rate of diffusion in different
places.

\section{Option G: Urban Environments}

\subsection{Q13: Fastest Growing Cities \score{10/10}}

\rowcolors{2}{TableStripe}{white}
\begin{tabularx}{\textwidth}{L{0.16\textwidth}L{0.12\textwidth}Y}
\toprule
\textbf{Part} & \textbf{Mark} & \textbf{Rationale} \\
\midrule
13(a)(i) & \score{1/1} & \term{Kinshasa} is the correct African city. \\
13(a)(ii) & \score{1/1} & \term{82} new people per hour for Shanghai is
correct. \\
13(b) & \score{2/2} & Increased greenhouse-gas emissions from more cars is a
valid environmental problem, with a direct development link to rapid population
growth. \\
13(c)(i) & \score{3/3} & Better job opportunities are a valid economic factor.
The student develops this through rural-urban migration and better-paying jobs,
then links it to urban population growth. \\
13(c)(ii) & \score{3/3} & The answer identifies a demographic structure: more
young people of working/childbearing age. It develops this through higher births
and links it to city growth. \\
\bottomrule
\end{tabularx}
\rowcolors{2}{}{}

\subsection{Q14(b): Centrifugal Population Movements \score{8/10}}

This response fits the 7--8 band and is awarded the upper end of that band. It
understands centrifugal population movement as movement away from the city
proper and applies it to Detroit. The example is relevant and includes useful
detail: suburbanization, white flight, automobile industry decentralization,
movement from the city to the outskirts, population decline from about
\term{1.7 million} to about \term{670,000}, loss of assembly plants, lower tax
revenue, harder infrastructure upkeep, poverty, crime and social tension.

The answer also addresses \term{to what extent} by recognizing benefits as well
as costs. It suggests that out-migration can reduce competition for housing or
resources for those who remain, possibly lowering prices and easing pressure on
infrastructure. This gives the answer some balance and a clear final judgment:
centrifugal movement generally burdens residential areas, but it may have some
limited benefits.

The response does not reach \score{9--10} because the evidence is mainly from
one inner-city decline case. A top-band answer would also cover suburban
residential impacts, such as urban sprawl, commuter villages, loss of green
space and habitat, and changing social characteristics in suburban areas. It
would also distinguish suburbanization from counter-urbanization and consider
stakeholders such as planners, developers, residents and environmental groups.

\section{Main Targets}

\begin{itemize}
  \item In 10-mark answers, cover a wider range of geographic factors rather
  than relying on one main pathway.
  \item Use precise case-study evidence and avoid historical shortcuts that
  weaken credibility.
  \item Keep the evaluative conclusion, but make it more explicit: state which
  factor or impact is most important, where, and why.
  \item For short-answer questions, name the exact category being credited, for
  example \term{human/manual energy} rather than simply \term{energy}.
\end{itemize}

\end{document}
